\documentclass[10pt,a4paper]{amsart}
\usepackage{amsmath,amssymb,amsthm,graphicx}
\usepackage{verbatim}
\newenvironment{reference}{\comment}{\endcomment}
%\newenvironment{reference}{}{}
\newenvironment{slogan}{\comment}{\endcomment}
\newenvironment{history}{\comment}{\endcomment}

% For commutative diagrams we use Xy-pic
\usepackage[all]{xy}

% We use 2cell for 2-commutative diagrams.
\xyoption{2cell}
% Theorem environments.
%
\theoremstyle{plain}
\newtheorem{theorem}[subsection]{Theorem}
\newtheorem{proposition}[subsection]{Proposition}
\newtheorem{lemma}[subsection]{Lemma}

\theoremstyle{definition}
\newtheorem{definition}[subsection]{Definition}
\newtheorem{example}[subsection]{Example}
\newtheorem{exercise}[subsection]{Exercise}
\newtheorem{situation}[subsection]{Situation}

\theoremstyle{remark}
\newtheorem{remark}[subsection]{Remark}
\newtheorem{remarks}[subsection]{Remarks}

\numberwithin{equation}{subsection}

% Macros
%
\def\lim{\mathop{\mathrm{lim}}\nolimits}
\def\colim{\mathop{\mathrm{colim}}\nolimits}
\def\Spec{\mathop{\mathrm{Spec}}}
\def\Hom{\mathop{\mathrm{Hom}}\nolimits}
\def\Ext{\mathop{\mathrm{Ext}}\nolimits}
\def\SheafHom{\mathop{\mathcal{H}\!\mathit{om}}\nolimits}
\def\SheafExt{\mathop{\mathcal{E}\!\mathit{xt}}\nolimits}
\def\Sch{\mathit{Sch}}
\def\Mor{\mathop{\mathrm{Mor}}\nolimits}
\def\Ob{\mathop{\mathrm{Ob}}\nolimits}
\def\Sh{\mathop{\mathit{Sh}}\nolimits}
\def\NL{\mathop{N\!L}\nolimits}
\def\CH{\mathop{\mathrm{CH}}\nolimits}
\def\proetale{{pro\text{-}\acute{e}tale}}
\def\etale{{\acute{e}tale}}
\def\QCoh{\mathit{QCoh}}
\def\Ker{\mathop{\mathrm{Ker}}}
\def\Im{\mathop{\mathrm{Im}}}
\def\Coker{\mathop{\mathrm{Coker}}}
\def\Coim{\mathop{\mathrm{Coim}}}

% Boxtimes
%
\DeclareMathSymbol{\boxtimes}{\mathbin}{AMSa}{"02}

%
% Macros for moduli stacks/spaces
%
\def\QCohstack{\mathcal{QC}\!\mathit{oh}}
\def\Cohstack{\mathcal{C}\!\mathit{oh}}
\def\Spacesstack{\mathcal{S}\!\mathit{paces}}
\def\Quotfunctor{\mathrm{Quot}}
\def\Hilbfunctor{\mathrm{Hilb}}
\def\Curvesstack{\mathcal{C}\!\mathit{urves}}
\def\Polarizedstack{\mathcal{P}\!\mathit{olarized}}
\def\Complexesstack{\mathcal{C}\!\mathit{omplexes}}
% \Pic is the operator that assigns to X its picard group, usage \Pic(X)
% \Picardstack_{X/B} denotes the Picard stack of X over B
% \Picardfunctor_{X/B} denotes the Picard functor of X over B
\def\Pic{\mathop{\mathrm{Pic}}\nolimits}
\def\Picardstack{\mathcal{P}\!\mathit{ic}}
\def\Picardfunctor{\mathrm{Pic}}
\def\Deformationcategory{\mathcal{D}\!\mathit{ef}}
\usepackage{cleveref}
\crefname{originallemma}{original lemma}{original lemmas}
\Crefname{originallemma}{Original lemma}{Original lemmas}
\makeatletter
\newlabel{more-algebra-lemma-regular-local-UFD}{{0AG0}{source}{}{sourcepointer-1}{}}
\newlabel{more-algebra-lemma-regular-local-UFD@cref}{{[originallemma][0][]0AG0}{[1][0][]source}{}{}{}}
\makeatother
\title{Ideals of a regular local ring of dimension 2 with reduced or depth-one quotient are principal or trivial}
\author{The Stacks Project Authors}
\date{stacks-project:curves.tex:local-snapshot-20261012; GNU FDL 1.2 or later}
\begin{document}
\maketitle
\noindent\textit{Editorial note.} All delivered mathematical text is the author's own, copied from the pinned Stacks Project chapter file curves.tex; no statement or proof step was rewritten, reordered or removed.\par\medskip
\noindent\textit{Editorial note.} The counted proof is the author's complete argument: the gcd factorization, the height-one and dimension dichotomy, the radical-case arguments, and the depth-one annihilator argument are all retained verbatim.\par\medskip
\noindent\textit{Editorial note.} The label lemma-reduced-quotient-regular-ring-dim-2 names the counted statement, and the proof environment immediately following it in the source is delivered verbatim as its proof.\par\medskip
\noindent\textit{Editorial note.} The counted statement is carried by the author's own lemma environment in the chapter source.\par\medskip
\noindent\textit{Allowed standard prerequisites (editorial).}\begin{itemize}
\item Regular local rings are UFDs. More on Algebra, Lemma more-algebra-lemma-regular-local-UFD. Provides the gcd factorization f\_i = f g\_i with gcd(g\_1, ..., g\_r) = 1 underlying the whole argument.
\item Dimension theory in a two-dimensional local ring. Standard dimension theory of Noetherian local rings (height-one primes and the maximal ideal as the only point of V(g) in dimension 2). Rules out height-one primes over the g\_i, leaving either (g\_i) = A or V(g\_i) = \{m\}.
\end{itemize}\medskip
\ifdefined\hypertarget\hypertarget{sourcepointer-1}{}\fi
\noindent\textit{Source reference.} Original lemma 0AG0: sources/stacks-chow/tags.txt, source lines 13178--13178. Prerequisite pointer invoked at proof line(s) 4047; it names a standard fact of a sibling chapter that the counted argument uses, and every step of the counted proof built on it is retained verbatim in the delivered body.\par\medskip
\begin{lemma}
\label{lemma-reduced-quotient-regular-ring-dim-2}
Let $(A, \mathfrak m)$ be a regular local ring of dimension $2$.
Let $I \subset \mathfrak m$ be an ideal.
\begin{enumerate}
\item If $A/I$ is reduced, then $I = (0)$, $I = \mathfrak m$, or
$I = (f)$ for some nonzero $f \in \mathfrak m$.
\item If $A/I$ has depth $1$, then $I = (f)$ for some nonzero
$f \in \mathfrak m$. 
\end{enumerate}
\end{lemma}

\begin{proof}
Assume $I \not = 0$. Write $I = (f_1, \ldots, f_r)$. As $A$ is a UFD
(More on Algebra, Lemma \ref{more-algebra-lemma-regular-local-UFD})
we can write $f_i = fg_i$ where $f$ is the gcd of $f_1, \ldots, f_r$.
Thus the gcd of $g_1, \ldots, g_r$ is $1$ which means that
there is no height $1$ prime ideal over $g_1, \ldots, g_r$.
Then either $(g_1, \ldots, g_r) = A$ which implies $I = (f)$ or
if not, then $\dim(A) = 2$ implies that
$V(g_1, \ldots, g_r) = \{\mathfrak m\}$, i.e.,
$\mathfrak m = \sqrt{(g_1, \ldots, g_r)}$.

\medskip\noindent
Assume $A/I$ reduced, i.e., $I$ radical. If $f$ is a unit, then since $I$
is radical we see that $I = \mathfrak m$. If $f \in \mathfrak m$, then we
see that $f^n$ maps to zero in $A/I$. Hence $f \in I$ by reducedness and
we conclude $I = (f)$.

\medskip\noindent
Assume $A/I$ has depth $1$. Then $\mathfrak m$ is not an associated
prime of $A/I$. Since the class of $f$ modulo $I$ is annihilated
by $g_1, \ldots, g_r$, this implies that the class of $f$ is zero
in $A/I$. Thus $I = (f)$ as desired.
\end{proof}

\end{document}
